% ABSTRACT (English)
This project aims to design and develop a voice-controlled ship navigation game as a learning tool to enhance communication skills and reduce foreign language anxiety among learners. The system integrates a simulation game format with speech recognition technology and consists of the following main systems: the speech recognition and command processing system, the ship physics and control system, the AI crew system, the obstacle generation system, the user interface, the language support system, the story mode, and the research logging system. The player takes the role of a ship captain who must communicate with and issue standard marine orders to an AI crew, namely the chief engineer, who runs the engine, and the helmsman and lookout, who steers the ship and watches for hazards. The game is developed with the Unity engine and C\#. English voice commands are transcribed on the player's computer by the Whisper model via whisper.cpp; unclear commands are interpreted with GPT-4o mini. The AI crew system simulates communication breakdowns: if the player's pronunciation is unclear or a command is out of context, the crew asks the player to confirm or repeat the order. The game is therefore designed to give players practice in speaking fluently and solving problems under time pressure. At this progress stage, a working prototype has been built, all 363 of 363 unit tests pass, and the developer's own system tests have been completed; the user evaluation is planned for the next phase. The system's performance will be measured from in-game logs, and the speaking anxiety of 10--20 participants will be compared before and after play using a 12-item scale adapted from the Foreign Language Classroom Anxiety Scale (FLCAS), together with a game experience questionnaire.
